\documentclass[12pt]{article}

% This is the preamble, load any packages you're going to use here
\usepackage{physics} % provides lots of nice features and commands often used in physics, it also loads some other packages (like AMSmath)
\usepackage{siunitx} % typesets numbers with units very nicely
\usepackage{enumerate} % allows us to customize our lists



\begin{document}

\title{Resonance Tube Lab Report}
\author{A.~Student, Another Student, Last Student}
\date{\today}

\maketitle

\begin{abstract}
	Eventually we will learn to write abstracts for our lab reports, but not today, please delete this abstract.
\end{abstract}

\section{Section Name}
	It is helpful to label the different sections of your report.
    Please choose names that make sense.  You may number the sections, like in this example or not (then put an asterisk after the section command, e.g.~\verb+\section*+). 
    
    
\section{Experimental Setup}
	This is a great section to explain the choices your group made to measure the resonance frequencies.  Use your answers to questions 1--4 to discuss these choices.
    
\section{Data}
	Please don't include any data tables in your \LaTeX\, write up.  Making tables in \LaTeX\, is very boring, although there are programs to convert your Excel file to \LaTeX\, form!  We'll worry about that later.  Simply submit your Excel or Google Sheets workbook via Sapling, sharing on Google drive, or email.
    

\section{Analysis}
	Here you discuss your observations and results.
	
    \subsection{One Part}
		The subsection command let's you further divide your sections up.  Comment on your observations and results here.
        
        You should write a few sentences about your answers to questions 6--8.
        
	\subsection{Second Part}
        Here's a good spot to discuss your plot (again, don't include the plot, we'll learn how to do this later).  Use your answer to question 7.
       
	\subsection{When to Make New Sections}
		As a general rule of thumb, don't make smaller sections, subsections, etc, unless there are at least two of them at that level (just like with lists).
	


\end{document}